\section{Conclusion}
\label{sec:conclusion}
Seven public phishing URL corpora were audited for presentation artefacts, sampling artefacts and registered-domain leakage. On PhiUSIIL a scheme-type and \texttt{www.} pattern separates the classes almost perfectly, also when both classes are restricted to homepages, and on the PT-Alexa construction the scheme flag alone does. Kaggle-MU mixes sources so that thousands of benign-looking pages carry the phishing label, and its formatting dependence cannot be separated from that mixing. Canonicalising the strings lowers a character n-gram detector by \DRawVsCanonPhiUSIIL{} MCC on PhiUSIIL, \DRawVsCanonKaggleMU{} on Kaggle-MU and \DRawVsCanonPTAlexa{} on PT-Alexa, and a character CNN loses \CNNDropPhiUSIILRandom{} and \CNNDropKaggleMURandom{} on the first two and \CNNDropPTAlexaRandom{} on PT-Alexa. On the other four corpora the formatting signal exists, but removing it costs the n-gram model less than 0.01 MCC on average and the CNN at most \CNNDropHannousseRandom{} under the random split. Grouping by registered domain costs the full-URL models up to \DLexicalRandomVsDomainKaggleMU{} MCC for the n-gram and lexical models and \CNNGroupCostISCXURL{} for the CNN, while it removes most of the score of a host-only model on ISCX-URL2016. The public files also contain duplicates, contained corpora and cross-file label conflicts that inflate or confuse results before any model is trained.

The practical message is modest. Raw-string scores on these corpora should not be compared across papers without the checks in Table~\ref{tab:checklist}, and Kaggle-MU should not serve as an independent benchmark. Canonical input helps transfer mainly where corpora use opposite conventions, and a threshold set for a low false-positive rate on one corpus does not hold that rate on another, so a detector moved to a new corpus would produce mostly false alarms at realistic prevalence until its threshold is set on the new traffic. Feature attribution ranks the corpora by how much a fitted detector uses formatting, including the two template corpora and Kaggle-MU, but it does not show whether a detector retrained without formatting would lose accuracy, so it can direct an audit but not replace the canonical intervention. A presentation probe, canonical and domain-grouped results and a purged transfer test take minutes with the audit code. Future work should repeat the audit on corpora in which both classes pass through one collection pipeline, such as telemetry-based collections, which would supply the positive control this study lacks, and should use collection dates instead of publication years.
